\chapter{Design Based Identification}
\label{chap:design}

\section{Instrumental variables}
\label{sec:design-iv}

When ignorability fails and the failure cannot be bounded tightly enough to be
useful, identification must come from the design\index{design based identification} rather than from adjustment.
The instrumental variable\index{instrumental variable} strategy exploits a
source of variation in treatment that is unrelated to the outcome except through
treatment.

\begin{assumption}[Instrument conditions]
\label{ass:iv}
A variable $Z$ is an instrument for the effect of $A$ on $Y$ if
$Z \indep \left( Y(a,z), A(z) \right)$ for all $a, z$, if
$\Prob(A(1) = 1) \neq \Prob(A(0) = 1)$, and if $Y(a, z) = Y(a)$ for all $z$.
\end{assumption}

The three conditions are independence, relevance, and
exclusion\index{exclusion restriction}. Only relevance
is testable. Exclusion is the substantive assertion, and in
Figure~\ref{fig:confounding-dag} it is the claim that the caseworker referral
$Z$ has no edge into earnings.

\begin{figure}[t]
\centering
\begin{tikzpicture}[
  >={Stealth[length=2.2mm]},
  obs/.style={draw,circle,minimum size=9mm,inner sep=0pt,font=\small},
  unobs/.style={draw,circle,dashed,minimum size=9mm,inner sep=0pt,font=\small},
  every edge/.style={draw,->,thick}
]
  \node[obs] (Z) {$Z$};
  \node[obs,right=20mm of Z] (A) {$A$};
  \node[obs,right=20mm of A] (Y) {$Y$};
  \node[unobs] (U) at ($(A)!0.5!(Y)+(0,-15mm)$) {$U$};

  \draw (Z) edge (A);
  \draw (A) edge (Y);
  \draw (U) edge (A);
  \draw (U) edge (Y);
  \draw[->,thick,dashed,red!60!black] (Z) to[bend left=38] node[above,font=\footnotesize] {excluded by assumption} (Y);

  \node[font=\footnotesize,below=1mm of Z,align=center] {referral};
  \node[font=\footnotesize,below=1mm of A,align=center] {enrolment};
  \node[font=\footnotesize,below=1mm of Y,align=center] {earnings};
  \node[font=\footnotesize,below=1mm of U,align=center] {motivation};
\end{tikzpicture}
\caption[Instrumental variable structure]{Instrumental variable structure for
the retraining study. The dashed arc is the edge that
Assumption~\ref{ass:iv} forbids. Identification rests entirely on its absence,
and no feature of the observed data distinguishes a world in which it is present
from one in which it is not.}
\label{fig:iv-dag}
\end{figure}

\begin{theorem}[Local average treatment effect]
\label{thm:late}
\index{local average treatment effect}
Under Assumption~\ref{ass:iv} and monotonicity, $A(1) \geq A(0)$ almost surely,
\begin{equation}
  \LATE \equiv \E\left[ Y(1) - Y(0) \,\middle|\, A(1) > A(0) \right]
    = \frac{\E[Y \mid Z=1] - \E[Y \mid Z=0]}{\E[A \mid Z=1] - \E[A \mid Z=0]}.
  \label{eq:late}
\end{equation}
\end{theorem}

\begin{proof}
Monotonicity\index{monotonicity} partitions the population into compliers with $A(1) > A(0)$, always
takers with $A(1) = A(0) = 1$, and never takers with $A(1) = A(0) = 0$, and
excludes defiers. By exclusion, $Y(a,z) = Y(a)$, so for always takers and never
takers the observed outcome does not depend on $Z$. By independence, the
subgroup proportions are the same in both arms of $Z$. Hence
$\E[Y \mid Z=1] - \E[Y \mid Z=0]$ equals the complier proportion times
$\E[Y(1) - Y(0) \mid \text{complier}]$, since the contributions of the other two
groups cancel. The denominator of~\eqref{eq:late} equals the complier proportion
by the same partition. Dividing gives the claim, and the denominator is nonzero
by relevance.
\end{proof}

\section{The estimand is not chosen}
\label{sec:design-late}

Theorem~\ref{thm:late} identifies an effect for compliers, a subgroup that
cannot be identified at the unit level and whose composition depends on the
instrument. Two valid instruments for the same treatment identify two different
quantities, and a disagreement between them is not evidence that one is invalid.

\begin{table}[t]
\centering
\caption[Instrument diagnostics and estimates]{Instrument diagnostics and
estimates for three candidate instruments in the retraining study. The complier
share is the first stage contrast of~\eqref{eq:late}. The final column reports
the change in the point estimate when the instrument is interacted with region,
which is a rough check on the exclusion restriction.}
\label{tab:iv-comparison}
\begin{tabular}{lrrrr}
\toprule
Instrument & First stage $F$ & Complier share & $\widehat{\LATE}$ & Interaction shift \\
\midrule
Caseworker referral        & 148.2 & 0.41 & 1604 & 84  \\
Distance to training site  & 31.7  & 0.19 & 2210 & 391 \\
Office opening month       & 8.4   & 0.06 & 4870 & 2914 \\
Referral and distance      & 91.5  & 0.47 & 1712 & 156 \\
\bottomrule
\end{tabular}
\end{table}

Table~\ref{tab:iv-comparison} illustrates the weak instrument problem\index{weak instruments} in a form
worth recognising. The office opening month has a first stage $F$ statistic of
8.4, below the conventional threshold of 10, and its point estimate is more than
three times the others. The naive reading is that opening month identifies a
different complier population. The correct reading is that a first stage this
weak amplifies any violation of exclusion by roughly the inverse of the complier
share, so a small direct effect of opening month on earnings would produce
exactly this pattern \citep{nakamura2015weakiv}. The interaction shift of 2914 in the final column is
larger than most of the point estimates in the table and is a straightforward
signal that the specification is not credible.

\section{Difference in differences}
\label{sec:design-did}

Where treatment is adopted at different times by different units, and where the
untreated trajectory of the treated units would have paralleled that of the
untreated, the effect is identified by a double difference.

\begin{assumption}[Parallel trends]
\label{ass:parallel}
\index{parallel trends}
$\E\left[ Y_{t}(0) - Y_{t-1}(0) \mid A = 1 \right]
 = \E\left[ Y_{t}(0) - Y_{t-1}(0) \mid A = 0 \right]$.
\end{assumption}

\begin{proposition}[Difference in differences]
\label{prop:did}
Under Assumptions~\ref{ass:consistency} and~\ref{ass:parallel},
\begin{equation}
  \ATT = \left( \E[Y_t \mid A=1] - \E[Y_{t-1} \mid A=1] \right)
       - \left( \E[Y_t \mid A=0] - \E[Y_{t-1} \mid A=0] \right).
  \label{eq:did}
\end{equation}
\end{proposition}

\begin{proof}
Add and subtract $\E[Y_t(0) \mid A=1]$ in the definition of the ATT. The
untreated period value satisfies $Y_{t-1} = Y_{t-1}(0)$ for both groups by
consistency, since treatment has not yet occurred. Assumption~\ref{ass:parallel}
replaces the unobserved $\E[Y_t(0) \mid A=1]$ by
$\E[Y_{t-1} \mid A=1] + \left( \E[Y_t \mid A=0] - \E[Y_{t-1} \mid A=0] \right)$.
Substituting and cancelling gives~\eqref{eq:did}.
\end{proof}

Assumption~\ref{ass:parallel} is not invariant to transformation of the outcome.
Trends that are parallel in levels are not parallel in logarithms unless the
groups have equal baseline means, so the scale of the outcome is part of the
assumption rather than a reporting choice. This is the single most common
unstated decision in applied difference in differences work
\citep{delacroix2021didscale}.

\section{The graphical view of design}
\label{sec:design-dag-view}

The strategies of this chapter and of Chapter~\ref{chap:graphical} are the same
mathematics under different labels. An instrument is a vertex with an edge into
$A$, no edge into $Y$, and no unblocked backdoor path to $Y$. Parallel trends is
the assertion that the confounder is additive and time invariant, which in a
diagram over $(Y_{t-1}, Y_t, A, U)$ means that $U$ has edges of equal weight
into both outcome vertices.

Stating the design assumption graphically is worth the effort because it makes
the failure modes visible. A time varying confounder breaks parallel trends by
giving $U$ unequal edge weights. An instrument that affects the outcome through
a second channel is an added edge. Neither failure is detectable in the data,
and both are easy to state, argue about, and if necessary bound with the
apparatus of Chapter~\ref{chap:sensitivity}.

\section{Interference}
\label{sec:design-interference}

Assumption~\ref{ass:sutva} has been maintained throughout. When it fails, the
potential outcome of unit $i$ is a function of the full assignment vector, and
the number of potential outcomes per unit grows exponentially in the sample
size. Progress requires restricting that dependence.

The most useful restriction is partial interference\index{interference!partial}: units fall into groups, and
spillover occurs within groups and not between them. Under partial interference
the direct effect and the spillover effect are separately identified by a
two stage randomisation, and in observational settings by a group level
propensity score \citep{halloran2017interference}. In the retraining study the natural group is the local labour
market office, and a displacement effect within an office would violate
Assumption~\ref{ass:sutva} in the direction that biases the estimate upward: a
trained worker who takes a job that an untrained worker would otherwise have
held raises the treated outcome and lowers the control outcome at once. Nothing
in Chapters~\ref{chap:propensity} through~\ref{chap:sensitivity} detects this,
and it is the most serious unaddressed threat to the worked example.
